\section{Conclusion}
A small open-weight model can learn to reproduce a product's decision policy on its benchmark from 820 synthetic examples on a free GPU. The learning is large, visible in every task and consistent across three base checkpoints: fine-tuning moved Hermes from 0.4755 to 0.8051 on a benchmark whose answers are computed by the policy, with a 95\% interval of 0.3068 to 0.3521 on the difference, replicating the pilot, and it moved Logos v0.0.1 by a similar amount. A stronger base alone did not produce a measurably better model; the Reasoning checkpoint of that base, trained on data that adds policy-derived reasoning traces, did, by $+0.0409$ on the answerable items with an interval that excludes zero but starts close to it. What the models learned is specific: they order and label the way the policy does, they decline by the kind of question rather than by the evidence, and they never produce the held-out format. Each model is one run, the second comparison confounds three changes, and the benchmark measures transfer across format within the scenarios the policy was written for. The contribution this paper claims is therefore not a method but an account: a pre-registered, fully logged record of what a student can measure about fine-tuning with the compute a student has, in which the per-task regressions, the inconclusive result and the measurably better one are reported to the same standard.
